\section{The framework}
\label{sec:framework}

\Cref{fig:framework} shows the three levels of the framework. The \emph{object level} improves a policy; the \emph{gate} decides what is kept; the \emph{meta level} improves the procedure the object level runs. Every control used in our experiments is a flag in the same engine, so a control differs from the arm it is compared with in exactly one place.

\subsection{Object level}
A \emph{program} $\pi$ is a small JSON document with four sections: \texttt{curriculum} (how many scenes to practise and how to allocate them across tasks), \texttt{solver} (samples per scene and sampling temperature), \texttt{curation} (which attempts become training items) and \texttt{update} (the operator and its hyperparameters). One \emph{inner step} of the loop, run from the current policy $\theta$:
\begin{enumerate}[nosep]
  \item \textbf{Propose.} Draw fresh scenes. In LIBERO a scene is a freshly sampled object placement for one of the ten tasks (seeds disjoint from every evaluation set); in the code domain the model writes a short program and an input.
  \item \textbf{Attempt.} The policy tries each scene $K$ times at temperature $T$.
  \item \textbf{Verify.} An external oracle the policy cannot change labels each attempt: the simulator's success flag, or a sandboxed interpreter.
  \item \textbf{Curate.} \texttt{mixed\_groups} keeps every attempt of scenes whose success rate lies strictly inside a band (used by GRPO); \texttt{verified} keeps successful attempts only (used by behaviour cloning, BC).
  \item \textbf{Update.} A LoRA update of $\theta$ with GRPO~\cite{shao2024deepseekmath} on the group-centred advantages, or BC/SFT on the kept attempts.
\end{enumerate}
A \emph{round} is three inner steps from the same parent, producing a \emph{candidate} version. The next round starts from the retained version, so practice is always generated by the latest policy: this is the \emph{recursion}.

\subsection{Gate, audit and test}
Three disjoint evaluation sets are drawn once per run. The \emph{gate} (20 placements per task, 200 episodes, greedy decoding, environment re-seeded per episode so that versions are compared on identical placements) decides: a candidate is accepted iff its gate success is at least the parent's. In the fresh-paired mode used for all main runs a new gate set is drawn every round and the parent is re-evaluated on it, so acceptance is always a paired comparison on placements never seen before. The \emph{audit} (another 200) is recorded but never decides anything; it exposes overfitting to the gate. The \emph{test} is the official benchmark (the 50 official initial states per task, 500 episodes) and is sealed: no process in the loop can read it, and it is read once, after every decision in every run of a study has been made (\cref{sec:protocol}).

\subsection{Meta level}
The meta level versions the program $\pi$ and lets it be improved under the same gate. Each round the incumbent program and $c{-}1$ proposed programs are each trained for real from the same parent, and the gate picks the version to keep; a program is credited only when its child wins. Three proposers are implemented:
\begin{itemize}[nosep]
  \item \texttt{fixed}: no change (the meta level is off).
  \item \texttt{numeric}: random mutation of declared knobs, biased toward knobs whose past changes raised the gate.
  \item \texttt{llm}: a language model (Claude Sonnet) reads the current program, the edit space with a neutral description of what each option computes, and, for every past candidate, its program diff, gate score, acceptance and per-step statistics (sampled success, share of mixed scenes, item counts, first/last training loss). It returns a diagnosis and one new program.
\end{itemize}
The edit space comes in two sizes. The \emph{knob} space has ten hyperparameters (scenes, allocation, samples, temperature, curation rule and budget, operator, learning rate, steps, batch). The \emph{harness} space adds nine components of the training procedure itself: the ratio computation of the policy gradient (legacy or matched), the loss aggregation (token or episode), the BC likelihood temperature (unit or rollout), the advantage normalisation, the curation band, the per-scene cap, the allocation floor and the two clip bounds. Each of these corresponds to a change that, during development, had been made by hand after diagnosing round records; the LLM is never told which option is better.

\paragraph{Attribution guard.} The meta level never sees the audit or test, the policy weights, the scenes or the episodes. It only chooses values inside the declared space; every proposal passes the same validator as a numeric mutation; invalid proposals and failed calls fall back to numeric search and are counted. A \emph{compute cap} scales any proposal down so that no candidate does more rollouts or more optimiser work per step than the starting program. All practice data still comes from the policy's own attempts and all verdicts from the oracle, so the LLM is a source of procedure choices, not of data.

\subsection{Controls}
\label{sec:controls}
\begin{description}[nosep]
  \item[\texttt{nonrec}] every round's practice is generated by the \emph{initial} policy (the update still continues from the retained version). Isolates recursion.
  \item[\texttt{shuffled}] success labels are permuted within each scene before curation. Isolates the verifier.
  \item[\texttt{onebig}] one round with the summed rollout and update budget of all rounds. Isolates iteration against plain one-shot RL.
  \item[\texttt{full}, \texttt{llm}, \texttt{llm\_harness}, \texttt{llm\_harness\_cap}] the meta level on, with the numeric proposer, the LLM over the knob space, the LLM over the harness space, and the latter under the compute cap. All have two candidates per round, so \texttt{full} and the \texttt{llm} arms cost the same.
\end{description}

\subsection{Domains}
\textbf{LIBERO / OpenVLA-OFT.} The policy is OpenVLA-OFT~\cite{kim2025oft} (7B, action chunks of 8), started from a supervised fine-tune on \emph{one} demonstration per task~\cite{li2025simplevlarl}, so that there is room to improve (about 59\% on LIBERO-Goal, 19\% on LIBERO-10). Rollouts and GRPO/BC training use the SimpleVLA-RL stack with LoRA ($r{=}32$). \textbf{Code world models.} The model predicts the output of a Python program on an input (the CRUXEval-O format~\cite{gu2024cruxeval}); it writes its own programs, the sandbox executes them, and an attempt is verified iff its answer equals the interpreter's. Qwen2.5-Coder 1.5B and 7B~\cite{hui2024qwen25coder}, Qwen3-8B~\cite{yang2025qwen3} and Granite-3.3-8B~\cite{granite2024} are trained with LoRA SFT on verified traces.

\begin{figure*}[t]
  \centering
  \fbox{\parbox{0.96\textwidth}{\centering\vspace{14mm}\textit{[Figure: the three levels. Left: object loop (propose $\to$ attempt $\to$ verify $\to$ curate $\to$ update) producing a candidate; middle: gate / audit / sealed test; right: meta level versioning the program, with the three proposers and the controls drawn as switches on the same engine.]}\vspace{14mm}}}
  \caption{\textbf{Framework.} One engine, every ingredient a switch. The object loop practises on fresh scenes with the latest policy; the gate accepts a candidate only if it does not regress on paired held-out placements; the meta level versions the training program and may let a numeric search or an LLM edit it inside a declared space. Controls (\texttt{nonrec}, \texttt{shuffled}, \texttt{onebig}) flip one switch each.}
  \label{fig:framework}
\end{figure*}
